\chapter{Implementation}
\label{chap:implementation}

We have identified OCaml Effect Handlers are similar to \lange. Therefore we can implement a contract system op top of OCaml and use preprocessor to rewrite the AST accordingly to the theoretical compiler. We first define a contract system in OCaml in the form of attributes. These annotations are defined similarly to the contract $\kappa$ in \langc.

\section{Examples}

In Listing \ref{code:example1} is the code annotated with contracts, including contracts for values, functions, higher-order function, and dependent contracts.

\begin{listing}
\inputminted{ocaml}{code/example1.ml}
\caption{Example contracts annotations in OCaml}
\label{code:example1}
\end{listing}


Annotations are built using an attribute \texttt{@contract} attached to the \texttt{let} expression. This attribute is only valid when the \texttt{let} expression is a value, or a function, and must be annotated with its types.

In the example, \texttt{x} is constrained to bigger than $10$. If the provided value of \texttt{x} breaks the contract, an exception is thrown as part of blame. The exception is checked lazily only when \texttt{x} is being used. \texttt{f1}, \texttt{f2}, \texttt{f3} are functions, and they are annotated with function contract, dependent contract, and higher-order dependent contract respectively. Their semantics are identical to the semantics higher-order contracts.

\section{Annotations}

OCaml attributes have a certain structure\footnote{Attributes must follow the
attribute-payload structure as defined in
\href{https://ocaml.org/manual/5.4/attributes.html}{https://ocaml.org/manual/5.4/attributes.html}}.
We chose to implement the annotations as a type expression in OCaml of the form
$\tau \rightarrow \tau$. This choice limits what we can put in as the
predicates. Therefore, we only allow predicates pre-defined as a function or
available/imported at the variable annotated.

We define the following annotations that are similar to different types of
contract: flat, function, dependent.

A flat contract is applied for a single value $\rho$, annotated using a single
name. This name is the function used to check the value.

A function is composabled by many flat contracts in the form $\rho \rightarrow
... \rightarrow \rho$. In OCaml a function can receive many arguments. These
arguments can be values or functions. Part of the OCaml type signature, which
we utilized for contract annotation, parenthesis can be used to separate
between the main function contract, and contracts for arguments that are
functions.

A dependent contract is strictly a function receiving a single argument $\rho
\rightarrow \rho\ dep$. This choice is made to keep the semantics similar to
\langc. Alternatively, we can choose to implement the dependent contract as
checking with all arguments, $\rho \rightarrow .. \rightarrow \rho\ dep$. For
this thesis, we keep the contract system simple and allow only for single
argument.

\section{Preprocessor Rewriting}

We implemented a rewritter as a preprocessor using
\texttt{ppxlib}\footnote{\href{https://ocaml-ppx.github.io/ppxlib/ppxlib/index.html}{https://ocaml-ppx.github.io/ppxlib/ppxlib/index.html}}.
The rewriter walks the AST of OCaml, finds and modifies valid
\texttt{@contract} annotated \texttt{let} expression.

The rewritter also updates the variable into a function receiving three extra
arguments \texttt{pos,neg,cloc} for blame labels. \texttt{pos} is the positive
label, which is the blamed party for a flat contract. \texttt{neg} is the
negative label. \texttt{cloc} is the contract location used for dependent
contract blame assignment. Blame is designed as an exception.

\begin{listing}
\begin{minted}{ocaml}
exception Blame of string
\end{minted}
\caption{Blame exception}
\label{code:blame_exception}
\end{listing}

In the theoretical compiler, a list of handlers is accumulated. A careful
reader may notice that most handlers have the same structure, with two
differences, the predicate, and the blame label. Therefore, as an optimization,
we design a single effect type accepting a \textit{tuple of three values}: the
blame label, the predicate function, and the value to check. We also define
unify handler for this effect which has the contract checking logic. These are
illustrated in Listing \ref{code:contract_effect}.

This is an optimization when contracts do not need shared state. We can augment
this with a state provided to the \texttt{contract\_checking} and apply
predicate with the state. At the moment, the lack the semantics about sharing
states between contracts, so the current implementation works only as a
demonstration.

\begin{listing}
\begin{minted}{ocaml}
type _ Effect.t +=
  | Flat : (string * ('a -> bool) * 'a) -> 'a Effect.t

let rec contract_checking : type a b.
  a Effect.t
  -> ((a, b) Effect.Deep.continuation -> b) option
  = fun eff ->
  let open Effect.Deep in
  match eff with
  | Flat (label, check, arg) ->
      Some (fun k ->
        if check arg
        then continue k arg
        else discontinue k (Blame label))
  | _ -> None
\end{minted}
\caption{Contract Effect}
\label{code:contract_effect}
\end{listing}

\subsection{Flat contract rewriting}

Flat contracts rewrite, in Listing \ref{code:rewrite-flat}, by converting the value to performing the contract effect \texttt{Flat} with the appropriate arguments. Blame labels are provided as argument to the orginal function such that when used, correct blame label is applied.

\begin{listing}
\inputminted{ocaml}{code/rewrite-flat.ml}
\caption{Flat contract rewrite}
\label{code:rewrite-flat}
\end{listing}

\subsection{Functional contract rewriting}

With flat contracts defined, a functional contract rewrite can be broken down
into contracts for its arguments and its return value, as illustrated in
Listing \ref{code:rewrite-function}. Therefore a function receiving 3
arguments, will be rewritten into 4 sub-contracts, 3 for its arguments, and 1
for its return value. An exception for this rule is when the function has a
function in one or more of its arguments. Then the rewrite must effectively
rewrite the argument contract into a function to rewrite recursively. When
rewriting into a function, argument names are generated. Blame labels for
arguments are ``flipped'' according to the \langc\ semantics.

\begin{listing}
\inputminted{ocaml}{code/rewrite-function.ml}
\caption{Functional contract rewrite}
\label{code:rewrite-function}
\end{listing}

\subsection{Dependent contract rewriting}

Dependent contract has a predicate that receives the argument value. Therefore,
we build up the correct contract predicate for the output contract. The value
must be passed in with a different labels set from the value used, which we
only need to create a new random variable. Not only that, because the dependent
argument runs inside a handler context, it lacks a handler because the contract
effect handler is pushed into the continuation. Therefore, we must wrap it into
a handler, as in the theoretical compiler model. Then we provide this to the
post-condition predicate as argument. The process is recursive, if the
dependent argument also has a dependent contract the dependent argument is also
wrapped inside a handler. An example rewrite is provided in Listing
\ref{code:rewrite-dep}.

\begin{listing}
\inputminted{ocaml}{code/rewrite-dep.ml}
\caption{Dependent contract rewrite}
\label{code:rewrite-dep}
\end{listing}

\subsection{Handling Contract Effects}

Using the defined \texttt{contract\_checking} user can wrap all possible code
path to handle all contract effects. User should be aware of functions with
contract effects called outside the handler. Therefore, we suggest wrapping the
whole program under \texttt{contract\_checking}. This is equivalent to Effect
Safety Compiler Theorem.

\subsection{Blame labels assignment}

All uses of the variable must now be appended with the three parameters
\texttt{pos}, \texttt{neg}, and \texttt{cloc}. It is by practice to use the
module name hosting the variable as the positive label, the module name of
usesite as the negative label, and the variable name itself as the contract
label. However, during the limitation of preprocessor having limited scoping
functionality, we hereby define a simpler strategy for blame assingment. There
are three cases that we have to resolve for blame assignment.

\begin{itemize}
\item Function's argument

This blame assignment follows the semantics of \langc, where positive and
negative labels are "swapped" with each other on function input. Because
function arguments are bounded inside a function, the generated code
applies the labels of the parent function in the correct order. This
process happens for dependent contract as well, where the dependent
argument is applied with the correct labels from its parent function.

\item Global variable

Once a protected functions/values are used in the user code, they must be
applied with the correct labels, else the blame report confuses the
programmer. There are 2 main cases here, (i) the function is defined used
in the same module, (ii) the function is defined and used in two different
modules. We examine the two cases and demonstrate how should the blame
labels be assigned.

In the same module scenario, the positive label is the module and the function.
The negative label is the code usage, which is the same module. Because our
implementation does not transform recursive functions, the only place in
the same module that can use a protected function/value is other function,
or in annonymous scope. If the protected function is used inside another
function, the negative label is the name of the module and the hosting
function. If it is being used in an annonymous scope, the negative label is
the name of the module.

In different modules scenario, the positive label remains the module and the
function. The negative label is the code usage. Similarly, we can use the
module name with the function name if unannoymous.

In both of these scenarios, the contract location label is the same as the
positive label.

\item Local variable

Because a contract can be annotated in any let expression, it is possible to
annotate a local variable with contract. In this scenario, the positive
label is the module, and the hosting function. The negative label is where
the variable is being used, which is highly likely the function itself.
Therefore all three labels are the same.

\end{itemize}
